\section{The Exact Period-Eight Spectral Edge}\label{sec:period-eight-edge}

The rational number in Section~\ref{sec:periodic-construction} is convenient for proving strict finite
counterexamples, but it is not the true band edge. We now prove Theorem~\ref{thm:exact-period-eight-edge}.

\subsection{The endpoint roots}

At \(c=2\), Equation~\eqref{eq:floquet-polynomial} becomes

\begin{equation}
\label{eq:edge-polynomial}
P(y,2)=y^{4}-16y^{3}+76y^{2}-96y+16.
\end{equation}

Translate \(y=x+4\). Then

\begin{equation}
\label{eq:edge-depressed-quartic}
P(x+4,2)=x^{4}-20x^{2}+80.
\end{equation}

Putting \(w=x^{2}\) gives \(w^{2}-20w+80=0\), whose roots are
\(10+-2\sqrt{5}\). The four real roots of \eqref{eq:edge-polynomial}, in increasing order, are

\begin{equation}
\label{eq:edge-four-roots}
\begin{aligned}
&4-\sqrt{10+2\sqrt{5}}, \\
&4-\sqrt{10-2\sqrt{5}}, \\
&4+\sqrt{10-2\sqrt{5}}, \\
&4+\sqrt{10+2\sqrt{5}}.
\end{aligned}
\end{equation}

Thus the largest endpoint root is

\begin{equation}
\label{eq:eta-definition}
\eta=4+\sqrt{10+2\sqrt{5}}.
\end{equation}

\subsection{Sharp positivity and uniqueness}

Set

\begin{equation}
\begin{aligned}
s&=\sqrt{10+2\sqrt{5}},       \eta=4+s, \\
u&=y-\eta,                   t=2-c.
\end{aligned}
\end{equation}

Substitution in \eqref{eq:floquet-polynomial} yields the identity

\begin{equation}
\label{eq:eta-positive-expansion}
\begin{aligned}
&P(\eta+u,2-t) \\
&=u^{4}+4s u^{3}+2u^{2}t+(40+12\sqrt{5})u^{2} \\
&\quad +4s ut+8\sqrt{5}s u+t^{2}+(4\sqrt{5}-3)t.
\end{aligned}
\end{equation}

All coefficients are positive. Hence for \(u,t\geq 0\), the right side is
nonnegative, and it vanishes only at \(u=t=0\). Since every unit-circle fiber
has \(c\leq 2\), Equation~\eqref{eq:eta-positive-expansion} excludes every squared eigenvalue above \(\eta\) and
permits equality only when \(c=2\).

At \(c=2\), Equation~\eqref{eq:edge-polynomial} has \(\eta\) as a root, so \(+-\sqrt{\eta}\) are eigenvalues
of \(H(1)\). Therefore

\begin{equation}
\label{eq:exact-infinite-edge}
\sup_{\lvert z\rvert=1} \rho(H(z))^{2}=\eta.
\end{equation}

On the unit circle, \(z+z^{-1}=2\) if and only if \(z=1\). This proves both the
value and the uniqueness assertion in Theorem~\ref{thm:exact-period-eight-edge}.

For later use we also record that \(\eta\) has minimal polynomial

\begin{equation}
\label{eq:eta-minimal-polynomial}
Y^{4}-16Y^{3}+76Y^{2}-96Y+16
\end{equation}

Indeed, translation \(Y=X+4\) reduces \eqref{eq:eta-minimal-polynomial} to
\(X^{4}-20X^{2}+80\). By Gauss's lemma, a rational factorization may be taken into
monic integral quadratics. Vanishing of the cubic and linear coefficients
reduces it either to \((X^{2}+a)(X^{2}+b)\), where
\(a+b=-20\) and \(ab=80\), or to
\((X^{2}+rX+s)(X^{2}-rX+s)\), where \(s^{2}=80\). The first case would make \(a,b\)
roots of \(T^{2}+20T+80\), whose discriminant 80 is not a square; the second has
no rational \(s\). Thus \eqref{eq:eta-minimal-polynomial} is irreducible over \(Q\). The number \(\eta\) lies in
\((\frac{1951}{250},\frac{1561}{200})\); in particular \(\eta<8\).

\subsection{Monotonicity of the top band}

Let \(r(c)\) be the largest real root of \(P(y,c)\). At \(c=-2\),

\begin{equation}
\label{eq:minus-one-fiber-factorization}
\begin{aligned}
P(y,-2)&=(y^{2}-12y+34)(y^{2}-4y+2), \\
r(-2)&=6+\sqrt{2}=:y_{0}.
\end{aligned}
\end{equation}

For \(-2<c\leq 2\), substitution gives

\begin{equation}
\label{eq:minus-one-fiber-comparison}
P(y_{0},c)=(c+2)(c+5-8\sqrt{2})<0.
\end{equation}

Since \(P(y,c)\) tends to positive infinity with \(y\), \eqref{eq:minus-one-fiber-comparison} implies
\(r(c)>y_{0}\) for \(c>-2\).

Moreover,

\begin{equation}
\label{eq:floquet-c-derivative}
P_c(y,c)=2(c-c_{0}(y)),       c_{0}(y)=y^{2}-8y+\frac{13}{2}.
\end{equation}

The function \(c_{0}\) is increasing for \(y\geq y_{0}>4\), and

\begin{equation}
c_{0}(y_{0})=-\frac{7}{2}+4\sqrt{2}>2.
\end{equation}

Thus \(P_c(y,c)<0\) throughout \(y\geq y_{0}\) and \(c\leq 2\). If
\(-2\leq c_{1}<c_{2}\leq 2\) and \(y_{1}=r(c_{1})\), then
\(P(y_{1},c_{2})<P(y_{1},c_{1})=0\); hence \(P(.,c_{2})\) has a root above \(y_{1}\).
Therefore \(r(c)\) is strictly increasing on \([-2,2]\).

\subsection{Finite holonomies}

For \(\alpha=+1\), the finite set \(z^L=1\) contains \(z=1\). By \eqref{eq:exact-infinite-edge},

\begin{equation}
\label{eq:positive-holonomy-edge}
\rho(A_{8L,+1})^{2}=\eta.
\end{equation}

For \(\alpha=-1\), the admissible parameters are
\(z_k=exp((2k+1)\pi i/L)\). Their largest \(c\) value is \(2\cos(\frac{\pi}{L})\), which is
strictly less than two. The monotonicity just proved gives

\begin{equation}
\label{eq:negative-holonomy-edge}
\rho(A_{8L,-1})^{2}=r(2\cos(\frac{\pi}{L}))<\eta.
\end{equation}

As \(L\) tends to infinity, \(2\cos(\frac{\pi}{L})\) tends to two, and continuity of the
eigenvalues of the Hermitian fiber gives convergence of \eqref{eq:negative-holonomy-edge} to \(\eta\).
Thus both holonomy sectors have the same infinite-volume edge, but only the
positive sector attains it at each finite size.
